\section{Graph Controlled Circuits}

\red{All circuits presented in this work are graph controlled circuits, which are a directed class of quantum circuits.
We here show, that their measurement distributions are equivalent with Bayesian networks.}

%Let us now introduce the most generic circuit construction schemes used in this work.

\subsection{Definition}

We build on the definition of uniformly controlled unitaries (see \cite{mottonen_transformation_2005}), and study their alignment along directed acyclic hypergraphs.
% Hypergraphs
We define graph-controlled circuits using hypergraphs, which nodes are associated with the qubits of the circuit.
The hyperedges are tuples of two subsets of the qubits, namely the incoming qubits representing control and the outgoing qubit representing the target of a unitary.
To each hyperedge we then associate a uniformly controlled unitary.

\begin{definition}[Graph-Controlled Circuit]
    Let $\graph=(\nodes,\edges)$ be a directed acyclic hypergraph, where each hyperedge has exactly one outgoing node and all nodes appear exactly once as outgoing nodes of an hyperedge.
    Then a by $\graph$ controlled circuit is a decoration of the edges $\edge=(\incomingnodes,\{\node\})\in\edges$ by uniformly controlled unitaries
    \begin{align*}
        \contunitaryofat{\edge}{\catvariableof{\node,\insymbol},\catvariableof{\node,\outsymbol},\catvariableof{\incomingnodes,\outsymbol}} \, .
    \end{align*}
\end{definition}

\begin{example}
    Each \activationCircuit{} is a graph-controlled circuit, where the corresponding hypergraph consists in $\nodes=\shortcatvariables\cup\{\avariable\}$ and a single edge $\edges=(\shortcatvariables,\{\avariable\})$.
\end{example}

% Orders of the qubits respecting directionality
Since the hypergraph $\graph$ to a graph-controlled circuit is acyclic, we find an order of the qubits, such that to each qubit appearing in the outgoing qubits of a hyperedge is ordered after the incoming qubits.
In such way, we have a unique definition of the graph-controlled circuit as a concatenation of the controlled unitaries respecting a directionality order.
When there are multiple orders respecting the directionality, then the concatenation respecting the orders all lead to an equivalent circuit.
% Should we show this using TN?

%%
\begin{example}[Generic State Representation]
    \label{exa:genericControlCircuit}
    In \cite{mottonen_transformation_2005} a preparation scheme for arbitrary states with real and positive amplitudes by graph-controlled circuits is presented.
    Therein, the hypergraph $\graph=(\nodes,\edges)$ is constructed as follows:
    \begin{itemize}
        \item Nodes $\nodes={[\catorder]}$
        \item Edges $\edges=\left\{([\catenumerator],\{\catenumerator\})\wcols\catenumeratorin\right\}$
    \end{itemize}
    It is possible to realize each controlled unitary by a sequence of single qubit rotations and C-NOTs (see \cite{mottonen_transformation_2005}), where the angles are computed using gray codes.
\end{example}


\subsection{Equivalence with Bayesian Networks}

%% Missing direction of Equivalence with BN
Above we have shown, how Bayesian Networks can be prepared by \activationCircuits{}, which are a class of graph-controlled circuits.
We now show, that the measurement distribution of any graph-controlled circuit is a Bayesian Network.
It then follows that the set of distributions prepared by graph-controlled circuits is exactly the set of Bayesian Networks.


This is closely connected to the chain decomposition of probability distributions, and their construction.
In \cite{low_quantum_2014} Bayesian Networks are prepared by graph-controlled circuits.
When in the graph $\graph$ each node appears at most once as outgoing node, it is also the hypergraph to a family of Bayesian Networks.
The measurement distributions of a state prepared by a $\graph$-controlled Circuit acting on disentangled initial states are exactly the Bayesian Networks with respect to $\graph$.

\begin{theorem}
    \label{the:graphControlledPreparesBN}
    Let $\graph$ be a directed acyclic hypergraph, such that node appears at most once as an outgoing node.
    The measurement distributions of the by $\graph$ controlled circuits acting on disentangled initial states are equal to the Bayesian Networks on $\graph$.
\end{theorem}

\begin{lemma}
    \label{lem:BNtoGCC}
    Let $\graph$ be a directed acyclic hypergraph, such that node appears at most once as an outgoing node.
    Then any Bayesian network on $\graph$ can be prepared by a $\graph$-controlled circuit with activation circuits of the conditional probability tensors.
\end{lemma}
\begin{proof}
    Let $\probwith$ be a Bayesian network on the graph $\graph$.
    Enumerate the nodes $\nodes$ of the $\graph$ by $[\atomorder]$, such that for each $\catenumeratorin$ we have $\parentsof{\catenumerator}\subset[\catenumerator]$.
    Then define a $\graph$-controlled circuit, by choosing for each $\catenumeratorin$ controlled unitaries which satisfy
    \begin{align*}
        \contunitaryofat{\catenumerator}{\ccatvariableof{\catenumerator}{\insymbol}=0,\catvariableof{\node,\outsymbol},\catvariableof{\parentsof{\catenumerator},\outsymbol}}
        =\sqrt{\condprobof{\catvariableof{\catenumerator}}{\catvariableof{\parentsof{\catenumerator}}}} \, .
    \end{align*}
    Here we specified only the action of the controlled unitary on the basis vector $\onehotmapofat{0}{\catvariableof{\catenumerator}}$, the action on $\onehotmapofat{1}{\catvariableof{\catenumerator}}$ can be chosen by an arbitrary orthogonal unit vector. % Maybe link here the activation circuit scheme, which is already used?
    For more explicit construction, see the \activationCircuits{}.
    Any such defined $\graph$-controlled circuit acting on the initial state $\bigotimes_{\catenumeratorin}\onehotmapofat{0}{\catvariableof{\catenumerator}}$ prepares a quantum state $\qstatewith$ with measurement distribution
    \begin{align*}
        \absof{\qstate}^2\left[\shortcatvariables\right] \, .
    \end{align*}
    Given arbitrary $\shortcatindicesin$ we have
    \begin{align*}
        \absof{\qstate}^2\left[\indexedshortcatvariables\right]
        = \prod_{\catenumeratorin} \condprobof{\indexedcatvariableof{\catenumerator}}{\indexedcatvariableof{\parentsof{\catenumerator}}}
        = \probat{\indexedshortcatvariables} \, .
    \end{align*}
    Here we used in the last equation, that $\probwith$ is a Bayesian network.
    Since the equivalence holds for any coordinate, this establishes the equivalence of the measurement distribution of $\qstatewith$ and $\probwith$.
\end{proof}

While we have already shown by \lemref{lem:BNtoGCC} that arbitrary Bayesian Networks can be prepared by graph-controlled circuits, we now show the converse, namely that any distribution prepared by graph-controlled circuits is a Bayesian Network.
To this end, we use the characterization of Bayesian Networks by the conditional independencies they encode (see \cite{koller_probabilistic_2009}).
To be more precise, we need to show that any node variable is conditionally independent of its non-descendants given its parents.

\begin{lemma}
    \label{lem:GCCtoBN}
    Let $(\graph,\contunitary)$ be a $\graph$-controlled circuit acting on a disentangled initial state and $\probat{\nodevariables}$ the corresponding measurement distribution.
    Then we have for each $\node\in\nodes$ the conditional independence
    \begin{align*}
        \condindependent{\catvariableof{\node}}{\catvariableof{\nondescendantsof{\node}}}{\catvariableof{\parentsof{\node}}} \, .
    \end{align*}
\end{lemma}
\begin{proof}
    We choose to a given $\node\in\nodes$ an enumeration $[\catorder]$ of the nodes, such that for each $\catenumeratorin$ we have $\parentsof{\catenumerator}\subset[\catenumerator]$ and for the enumerator $\seccatenumerator$ of $\node$ we further have $\nondescendantsof{\seccatenumerator}\subset[\seccatenumerator]$.
    Let $\probwith$ be the measurement distribution of the $\graph$-controlled circuit acting on a disentangled initial state $\bigotimes_{\catenumeratorin}\qstateofat{\catenumerator}{\catvariableof{\catenumerator}}$ and choose arbitrary $\catindexof{[\catenumerator]}$.
    We then have
    \begin{align*}
        \probat{\catvariableof{\seccatenumerator},\indexedcatvariableof{[\catenumerator]}}
        &= \contractionof{\left(\bigcup_{\catenumeratorin}\{\contunitaryof{\catenumerator},\contunitaryof{\catenumerator,\dagger},\qstateof{\catenumerator},\qstateof{\catenumerator,*}\right)\}
            \cup \left(\bigcup_{\catenumerator\in[\seccatenumerator]} \onehotmapofat{\catindexof{\catenumerator}}{\ccatvariableof{\catenumerator}{\outsymbol}}\right)
        }{
            \catvariableof{\seccatenumerator}
        } \\
        &= \contractionof{\bigcup_{\catenumerator\in[\seccatenumerator]}\{\contunitaryof{\catenumerator},\contunitaryof{\catenumerator,\dagger},\qstateof{\catenumerator},\qstateof{\catenumerator,*},\onehotmapofat{\catindexof{\catenumerator}}{\ccatvariableof{\catenumerator}{\outsymbol}}\}}{
            \catvariableof{\seccatenumerator}
        } \\
        & = \absof{\contractionof{\contunitaryofat{\seccatenumerator}{\catvariableof{\seccatenumerator,\insymbol},\catvariableof{\seccatenumerator,\outsymbol}},\indexedcatvariableof{\parentsof{\seccatenumerator},\outsymbol}}{\catvariableof{\seccatenumerator,\outsymbol}}}^2 \\
        & \quad \quad \cdot \prod_{\catenumerator\in[\seccatenumerator]}
        \left(\contraction{\qstateofat{\catenumerator}{\ccatvariableof{\catenumerator}{\insymbol}},\contunitaryofat{\catenumerator}{\ccatvariableof{\catenumerator}{\insymbol},\ccatvariableof{\catenumerator}{\outsymbol},\indexedcatvariableof{\parentsof{\catenumerator},\outsymbol}}}\right)^2
        %\contraction{\qstateof{\catenumerator,*},\onehotmapofat{\catindexof{\catenumerator}}{\catvariableof{\catenumerator}}}{} \, .
    \end{align*}
    Here we used in the second equation the unitarity of the controlled unitaries to $\catenumerator\notin[\seccatenumerator]$.
    Since the indices $\catindexof{[\seccatenumerator]/\parentsof{\seccatenumerator}} = \catindexof{\nondescendantsof{\seccatenumerator}}$ appear only in the constant term, we conclude
    \begin{align*}
        \condprobat{\catvariableof{\seccatenumerator}}{\catvariableof{[\catenumerator]}} = \condprobat{\catvariableof{\seccatenumerator}}{\catvariableof{\parentsof{\seccatenumerator}}} \otimes \onesat{\catvariableof{\nondescendantsof{\seccatenumerator}}}\, ,
    \end{align*}
    which establishes the conditional independence $\condindependent{\catvariableof{\node}}{\catvariableof{\nondescendantsof{\node}}}{\catvariableof{\parentsof{\node}}}$.
\end{proof}

\begin{proof}[Proof of \theref{the:graphControlledPreparesBN}]
    The theorem follows directly from \lemref{lem:BNtoGCC} and \lemref{lem:GCCtoBN}, using that Bayesian Networks are characterized by the conditional independence of each variable to its non-descendants given its parents.
\end{proof}

\begin{example}[Chain decomposition]
    The hypergraph in \exaref{exa:genericControlCircuit} corresponds with a family of Bayesian Networks subsuming any probability distribution.
    This can be verified based on the chain decomposition of a generic distribution.
\end{example}

%% Phase representation
Another question is, whether each quantum state, which measurement distribution is a Bayesian Network can be prepared by a $\graph$-controlled circuit.
This is not always the case, since the phase tensor does not influence the measurement distribution.
Any phase tensor, of a by $\graph$-controlled circuit prepared state has however a decomposition
\begin{align*}
    \phasecoreat{\shortcatvariables}
    = \sum_{\catenumeratorin} \phasecoreofat{\catenumerator}{\catvariableof{\catenumerator},\catvariableof{\parentsof{\catenumerator}}} \otimes \onesat{\catvariableof{[\catorder]/\{\{\catenumerator\}\cup\parentsof{\catenumerator}\}}} \, ,
\end{align*}
where the phase cores $\phasecoreof{\catenumerator}$ can be read of the controlled unitaries.
When there are phase cores which do not have such a decomposition, the corresponding states are not representable.
\cite{mottonen_transformation_2005} shows, that when the graph is chosen as in \exaref{exa:genericControlCircuit}, then also the phases can be represented.

